\subsection{局部环}

命题 1. --- 设 A 是非零环，r 是 A 的不可逆元素组成的集合。下列性质等价：

\begin{enumerate}
\def\labelenumi{(\roman{enumi})}
\item
  集合 \(\mathfrak{r}\) 是 A 的一个双边理想。
\item
  集合 \(\mathfrak{r}\) 对加法封闭。
\item
  环 A 有唯一的极大左理想。
\item
  对每个 \(a \in \mathbf{A}\) ，元素 \(a\) 与 \(1 - a\) 之一可逆。
\item
  对每个 \(a \in \mathbf{A}\) ，元素 \(a\) 与 \(1 - a\) 之一左可逆。
\end{enumerate}

蕴涵关系 (i) \(\Rightarrow\) (ii) 由理想的定义得出。由于 1 不属于 \(\mathfrak{r}\) ，我们有 (ii) \(\Rightarrow\) (iv)。

我们有 \(\mathfrak{r} \neq \mathrm{A}\) ，且集合 \(\mathfrak{r}\) 包含 \(\mathrm{A}\) 的每个不等于 \(\mathrm{A}\) 的左理想。若 \(\mathfrak{r}\) 是 \(\mathrm{A}\) 的一个左理想，则它因此是 \(\mathrm{A}\) 的唯一极大左理想。这就证明了 (i) 蕴涵 (iii)。

设 A 有唯一的极大左理想 m。取 \(b \in A - m\) 。左理想 Ab 不含于 A 的任何极大左理想，因此等于 A（I, §8, No.~6, p.~104, 定理 1），于是 b 左可逆。对每个 \(a \in A\) ，元素 a 与 1 - a 之一属于 A - m，因为 m 是一个不含 1 的理想。于是 (iii) 蕴涵 (v)。

设性质 (v) 成立。设 b 是 A 的一个左可逆元素。取 \(c \in A\) 使得 cb = 1。我们有 \((1 - bc)b = 0\) 且 \(b \neq 0\) ；因此 1 - bc 不是左可逆的。由性质 (v)，bc 左可逆，从而 c 更是左可逆的。于是 c 可逆；b 是它的逆，故 b 可逆。由此得出 (v) 蕴涵 (iv)。

剩下要证明 (iv) 蕴涵 (i)。设 (iv) 成立。则由下列断言 a) 至 d)，\(\mathfrak{r}\) 是 A 的一个双边理想：

\begin{enumerate}
\def\labelenumi{\alph{enumi})}
\item
  我们有 \(0 \in \mathfrak{r}\) ，因为环 A 非零。
\item
  A 的两个元素——一个属于 r，另一个属于 A-r——的乘积属于 r。
\item
  集合 \(\mathfrak{r}\) 对加法封闭。
\end{enumerate}

事实上，设 \(a\) 与 \(b\) 是 \(\mathfrak{r}\) 的元素，使得 \(s = a + b\) 可逆。由 b)，A 的元素 \(s^{-1}a\) 与 \(s^{-1}b\) 都属于 \(\mathfrak{r}\) ；由于 \(s^{-1}b = 1 - s^{-1}a\) ，这与 (iv) 成立的假设矛盾。

\begin{enumerate}
\def\labelenumi{\alph{enumi})}
\setcounter{enumi}{3}
\tightlist
\item
  集合 \(\mathfrak{r}\) 对乘法封闭。
\end{enumerate}

事实上，设 \(a\) 与 \(b\) 是 \(\mathfrak{r}\) 的两个元素。元素 \(a' = -a(1 - b)\) 属于 \(\mathfrak{r}\) （由 \(\mathrm{b}\) )），于是元素 \(ab\) ——它等于 \(a + a'\) ——属于 \(\mathfrak{r}\) （由 \(\mathrm{c}\)）；断言 d) 得证。

定义 1. --- 局部环是一个具有命题 1 中诸等价性质的非零环。

环 A 是局部环，当且仅当反环 \(A^{o}\) 是局部环。

若 A 是局部环，则 A 的不可逆元素组成的集合 r 是 A 的一个双边理想；它包含 A 的每个不等于 A 的左理想或右理想。因此环 \(A/r\) 是一个体，我们称之为 A 的剩余域。集合 r 是 A 的唯一极大左（相应地，右、双边）理想；我们简单地称 r 为 A 的极大理想。

例. --- 1) 每个体都是局部环。

\begin{enumerate}
\def\labelenumi{\arabic{enumi})}
\setcounter{enumi}{1}
\tightlist
\item
  设 A 是一个非零环，其中每个元素都可逆或幂零。则 A 是局部环。事实上，若 \(a \in A\) 不可逆，则由假设，存在整数 \(n \geqslant 0\) ，使得 \(a^{n+1} = 0\) ，而 1 - a 有逆元 \(1 + a + \cdots + a^{n}\) 。
\end{enumerate}

*3) 设 \(\mathrm{X}\) 是一个 \(\mathrm{C}^r\) 流形（VAR, R, 5.1.5），\(x\) 是 \(\mathrm{X}\) 的一点。设 \(\mathcal{O}_x\) 是在 \(x\) 处的 \(\mathrm{C}^r\) 函数（取值于标量体 \(\mathrm{K}\) ）的芽组成的环。则 \(\mathcal{O}_x\) 是交换局部环，其极大理想由在 \(x\) 处为零的函数的芽组成。

\begin{enumerate}
\def\labelenumi{\arabic{enumi})}
\setcounter{enumi}{3}
\item
  设 A 是交换局部环，\(B = A[[X_{i}]]_{i \in I}\) 是系数在 A 中的形式幂级数代数（III, §2, No.~11, p.~456）。由 IV, §4, No.~4, p.~30 的命题 6，环 B 是局部的，其极大理想由常数项属于 A 的极大理想的形式幂级数组成。特别地，若 A 是体，则 \(A[[X_{i}]]_{i \in I}\) 的极大理想由常数项为零的形式幂级数组成。
\item
  设 p 是素数。我们用 \(\mathbf{Z}_{(p)}\) 表示有理数体 Q 的由分式 a/b（其中 \(a \in Z\) 、 \(b \in Z\) 且 b 不被 \(p\) 整除）组成的子环*（参看 Comm. Alg., II, §2, No.~1, p.~60）*。则 \(\mathbf{Z}_{(p)}\) 是交换局部环，其极大理想为 \(p\mathbf{Z}_{(p)}\) 。p-adic 整数的环 \(Z_{p}\) （V, §12, No.~3, p.~96）是交换局部环，其极大理想为 \(pZ_{p}\) （VIII, p.~40, 习题 9）。
\item
  设 K 是特征 p \textgreater{} 0 的交换域，G 是一个 p-群（I, §6, No.~5, p.~76, 定义 9）。群 G 在 K 上的代数 K{[}G{]}（III, §2, No.~6, p.~446）是局部环；其极大理想是这样的元素 \((a_{g})_{g \in \mathrm{G}}\) ——K{[}G{]}中满足 \(\sum_{g \in G} a_{g} = 0\) 的元素——组成的集合（VIII, p.~41, 习题 10）。
\end{enumerate}

